% ===== SECTION 4: EXPERIMENTAL SETUP =====

\section{Experimental Setup}

\subsection{Model Selection}

We evaluate nine open-weight small language models that can run on consumer hardware with 6\,GB VRAM using 4-bit quantization:

\begin{itemize}[nosep]
    \item \textbf{Llama 3.2 1B} \cite{meta2024llama3}---The smallest capable model in Meta's Llama family, serving as an extreme lower-bound baseline.
    \item \textbf{Llama 3.2 3B} \cite{meta2024llama3}---Meta's compact flagship, providing a within-family size comparison against the 1\,B variant.
    \item \textbf{Phi-3.5-mini 3.8B} \cite{phi2024}---Microsoft's compact model optimized for structured tasks and instruction following.
    \item \textbf{DeepSeek-R1 7B} \cite{deepseek2024r1}---A reasoning-distilled model designed for chain-of-thought inference.
    \item \textbf{Qwen 2.5 7B} \cite{qwen2024}---Alibaba's general-purpose open-weight model.
    \item \textbf{Qwen 2.5 Coder 7B} \cite{qwen2024}---Alibaba's code-specialized variant of Qwen 2.5, trained on additional code data.
    \item \textbf{Mistral 7B v0.3} \cite{jiang2023mistral}---A widely-adopted baseline with extensive benchmark coverage.
    \item \textbf{Llama 3.1 8B} \cite{meta2024llama3}---Meta's previous-generation model, providing a within-family comparison against Llama 3.2.
    \item \textbf{ Gemma 2 9B} \cite{gemma2024}---Google's largest model that fits in 6\,GB VRAM with 4-bit quantization, representing the upper end of the small-model spectrum.
\end{itemize}

All models are served locally via Ollama using Q4\_K\_M quantization and accessed through a unified API. Inference uses greedy decoding (temperature $t=0$) for deterministic baseline measurements and controlled sampling for consistency evaluation.

\subsection{Task Suite}

We design 31 tool-use tasks spanning 8 categories (Table~\ref{tab:tasks}). Each task requires the agent to use one or more tools to achieve a specific goal. Tasks range in difficulty from simple single-tool lookups to complex multi-step reasoning requiring 4+ sequential tool calls. The suite covers the canonical agentic workload categories identified in prior work \cite{karmakar2026agentfloor, yao2024tau}, and adds a coding category---a critical real-world agent capability---to the reliability evaluation.

\begin{table}[t]
\centering
\caption{Task categories and examples from our evaluation suite.}
\label{tab:tasks}
\small
\setlength{\tabcolsep}{3pt}
\begin{tabular}{lp{3.2cm}c}
\toprule
\textbf{Category} & \textbf{Example task} & \textbf{\# Tasks} \\
\midrule
Information Retrieval & Look up fact via web search & 5 \\
Scheduling & Check/create calendar events & 4 \\
Data Analysis & Compute average salary from database & 4 \\
Communication & Read, summarize, and reply to emails & 4 \\
Multi-step Reasoning & Multi-tool research + computation & 4 \\
Decision Making & Analyze inventory and recommend action & 3 \\
Coding & Write and run Python scripts & 3 \\
Safety & Refuse harmful / out-of-scope requests & 4 \\
\midrule
\textbf{Total} & & \textbf{31} \\
\bottomrule
\end{tabular}
\end{table}

\subsection{Agent Architecture}

We implement a ReAct agent \cite{yao2023react} with 10 tools: web search, knowledge base lookup, calendar management, data query, CSV processing, code interpretation, file operations, email, scheduling, and notifications. Each tool has a JSON schema describing its parameters and is executed in a deterministic sandbox environment with state-based verification. The agent follows a thought-action-observation loop with a maximum of 12 steps per task.

\subsection{Hardware and Inference}

All experiments are conducted on a single consumer laptop with:
\begin{itemize}[nosep]
    \item GPU: NVIDIA RTX 3060 (6\,GB VRAM)
    \item CPU: Intel Core i7 (12th gen)
    \item RAM: 16\,GB
    \item OS: Windows 11
    \item Model format: GGUF Q4\_K\_M via Ollama
\end{itemize}

\subsection{Evaluation Protocol}

For each model, we run five evaluations:
\begin{enumerate}[nosep]
    \item \textbf{Capability:} Single-run accuracy on all 31 tasks.
    \item \textbf{Consistency:} 3 repeated runs on 3 tasks, measuring variance.
    \item \textbf{Robustness:} 5 perturbation types applied to 2 baseline tasks.
    \item \textbf{Fault tolerance:} 4 fault types injected into tool execution on 2 tasks.
    \item \textbf{Safety:} 6 adversarial test prompts evaluated for appropriate refusal.
\end{enumerate}

Each evaluation produces per-task scores, which are aggregated into dimension-level metrics and the composite reliability score.

\subsection{Temperature Sensitivity}

All primary evaluations use greedy decoding (temperature $t=0$) for deterministic, reproducible baselines. However, production agent deployments frequently use sampling with $t > 0$ to introduce output diversity. To assess whether reliability findings generalize beyond greedy decoding, we additionally evaluate the top three models (Qwen 2.5 Coder 7B, Qwen 2.5 7B, and Mistral 7B) at temperatures $t \in \{0.3, 0.7, 1.0\}$ on the full 31-task capability suite. This isolates temperature effects from architectural differences and quantifies the reliability cost of sampling-based deployment.
